\subsection{乘法代表元}

命题 6．——设 C 是一个 p-环，其剩余域 k 是完美的。假设 C 的长度有限为 n（分别为无限长度）。存在唯一同构 \(u:W_{n}(k)\to C\)（分别为 \(u:W(k)\to C\)），它经过取商诱导出 k 的恒等映射。

由于 \(W_{n}(k)\)（分别为 \(W(k)\)）是剩余域为 k 且长度为 n（分别为无限长度）的 p-环（第 1 号，例 2），又由于 \(\varnothing\) 是完美域 k 的 p-基，命题 6 是第 3 号命题 4 的特例。

定理 2．——设 A 是分离且完备的局部环，k 是其剩余域，π 是从 A 到 k 的典范同态。假设 k 是特征为 p 的完美域。

\begin{enumerate}
\def\labelenumi{\alph{enumi})}
\item
  存在唯一的环同态 \(u: \mathbf{W}(k) \to \mathbf{A}\)，使得有 \(\pi(u(\boldsymbol{a})) = a_0\)，其中 \(\boldsymbol{a} = (a_n)_{n \in \mathbb{N}}\) 属于 \(\mathbf{W}(k)\)。
\item
  当 W(k) 赋予 pW(k)-进拓扑时，同态 u 是连续的，且 u 的像是 A 的唯一科恩子环。
\end{enumerate}

根据第 2 号的定理 1，存在唯一的 A 的科恩子环；记为 C。设 u 是从 W(k) 到 A 的同态，使得有 \(\pi(u(\boldsymbol{a})) = a_0\)，对每个 \(\boldsymbol{a} = (a_n)_{n \in \mathbb{N}}\)（属于 W(k)）均如此；显然 u 的像是 A 的科恩子环，因而等于 C。u 的存在性和唯一性由命题 6 得出。C 的 pC-进拓扑由 A 的 \(m_{A}\)-进拓扑诱导（第 2 号，定理 1，b)），从而 u 连续。

关于 u 的直接构造，参见 p. 70，习题 6。

命题 7．——沿用定理 2 的假设与记号。存在唯一的 A 的乘法子集 S，使 π 将 S 双射到 k。A 的元素 a 属于 S 的充要条件是：对于每个 n ∈ N，存在 A 中的元素 \(a_{n}\)，使 \(a = a_{n}^{p^{n}}\)。集合 S 正是所有形如 \(u(x, 0, 0, \ldots)\) 的元素所成的集合。

先证明 S 的唯一性。设 S 是 A 的一个乘法子集，使 \(\pi\) 将 S 双射到 k。令 T 为 A 中这样的元素所成的集合：它们对每个 \(p^{n}\)（其中 \(n \in N\)）都是 p＾n 次幂。

\begin{enumerate}
\def\labelenumi{\alph{enumi})}
\item
  有 S ⊂ T：令 \(a \in S\) 且 \(n \in N\)；由于域 k 是完美的，存在 k 中的元素 \(x_{n}\)，使 \(x_{n}^{p^{n}} = \pi(a)\)；又由于有 \(\pi(S) = k\)，存在 S 中的元素 \(a_{n}\)，使 \(x_{n} = \pi(a_{n})\)。于是有 \(\pi(a_{n}^{p^{n}}) = \pi(a)\)，从而 \(a_{n}^{p^{n}} = a\)，因为 \(\pi\) 在 S 上的限制是单射。
\item
  \(\pi\) 在 T 上的限制是单射：设 \(a\) 和 \(b\) 是 T 的两个元素，满足 \(\pi(a) = \pi(b)\)。令 \(n \in \mathbf{N}\)；存在 A 中的两个元素 \(a_n\) 和 \(b_n\)，使 \(a = a_n^{p^n}\)，\(b = b_n^{p^n}\)。于是有 \(\pi(a_n)^{p^n} = \pi(b_n)^{p^n}\)，从而 \(\pi(a_n) = \pi(b_n)\)，即 \(a_n \equiv b_n\) 模 \(\mathfrak{m}_{\mathbf{A}}\)。根据 § 1 第 1 号的引理，有 \(a_n^{p^n} \equiv b_n^{p^n}\) 模 \(\mathfrak{m}_{\mathbf{A}}^{n+1}\)，即 \(a \equiv b\) 模 \(\mathfrak{m}_{\mathbf{A}}^{n+1}\)。由于 \(n\) 任意，有 \(a = b\)。
\end{enumerate}

上面的性质 a) 和 b)，连同公式 \(\pi(S) = k\)，蕴含关系 \(S = T\)，从而得到唯一性。

现在证明 S 的存在性。沿用定理 2 的记号，置 \(\varphi = u \circ \tau_k\)，即（§ 1，第 6 号）

\[
\varphi (x) = u (x, 0, 0, \dots)\tag{2}
\]

对于每个 \(x \in k\)。根据同上命题 4，有

\[
\varphi (1) = 1, \quad \varphi (x y) = \varphi (x) \varphi (y) \quad \text { for } \quad x, y \text {   in   } k.\tag{3}
\]

显然，映射 \(\pi \circ \varphi\) 是 \(k\) 的恒等映射。因此 \(\varphi\) 的像 S 满足命题 7 的条件。

S 中的元素通常称为 A 的乘法代表元（或 Teichmüller 代表元）。

注记．——1) 沿用前面的假设与记号。有

\[
\boldsymbol {a} = \sum_ {n = 0} ^ {\infty} p ^ {n} \tau_ {k} (a _ {n} ^ {p ^ {- n}}) \quad (\boldsymbol {a} = (a _ {n}) _ {n \in \mathbf {N}} \in \mathrm{W} (k))
\]

根据 § 1 第 8 号的命题 7。因此有

\[
u (\boldsymbol {a}) = \sum_ {n = 0} ^ {\infty} p ^ {n} \varphi (a _ {n} ^ {p ^ {- n}})\tag{4}
\]

对于每个 \(\boldsymbol{a} = (a_n)_{n \in \mathbf{N}}\)（属于 \(\mathrm{W}(k)\)），由于 \(u\) 连续（定理 2，b)），都有上式。根据公式（4），A 的唯一科恩子环由形如 \(\sum_{n=0}^{\infty} p^n s_n\) 的元素组成，其中 \(s_n \in S\) 对每个整数 \(n \geqslant 0\) 均成立。

\begin{enumerate}
\def\labelenumi{\arabic{enumi})}
\setcounter{enumi}{1}
\tightlist
\item
  设 A 是分离且完备的局部环，k 是其剩余域，且 \(\pi\) 是从 A 到 k 的典范同态。可以证明，存在 A 的一个乘法子集 S（一般不唯一），使 \(\pi\) 将 S 双射到 k（参见 p. 72，习题 11）。
\end{enumerate}

例．——1) 设 \(k\) 是特征为 \(p\) 的完美域。环 \(\mathbf{W}(k)\) 的乘法代表元是 Witt 向量 \(\tau(x) = (x, 0, 0, \ldots)\)，其中 \(x \in k\)。

\begin{enumerate}
\def\labelenumi{\arabic{enumi})}
\setcounter{enumi}{1}
\item
  设 A 是整的、分离且完备的局部环。假设 A 的剩余域 \(k\) 有限，含有 \(q = p^{f}\) 个元素，因而是特征为 \(p\) 的完美域。有 \(x^{q} = x\)，对每个 \(x \in k\) 均成立；从而有 \(s^{q} = s\)，对每个乘法代表元 \(s\) 均成立。由此，乘法代表元的集合由 0 以及 A 的分式域中的 \((q - 1)\) 次单位根组成。若 A 的分式域是局部紧的，则乘法代表元的存在性也可由 VI，§ 9，第 2 号，命题 3 得到（另参见 VI，§ 9，习题 5）。
\item
  更具体地，考虑 \(A = Z_{p}\) 的情形。此时乘法代表元是 0 以及 \((p - 1)\) 次单位根，它们位于分式域 \(Q_{p}\)（即 \(Z_{p}\) 的分式域）中。
\end{enumerate}

